\documentclass[a4paper,10pt]{scrartcl}

\usepackage[ngerman]{babel}
\usepackage[utf8]{inputenc}
\usepackage[T1]{fontenc}
\usepackage[babel,german=quotes]{csquotes}
\usepackage{hyperref}
\usepackage{geometry}
\usepackage{graphicx}
\usepackage{enumitem}
\usepackage{multicol}

\setlength\parindent{0pt}
\setitemize{itemsep=0pt}


\geometry{
	a4paper,
	left=20mm,
	top=20mm,
	bottom=15mm,
	footskip=5mm
}

\subject{Gruppe \textit{Nummer}}
\title{Benutzerhandbuch für das Thema \glqq\textit{Titel}\grqq}
\subtitle{Betreuer: \textit{Name}}
\author{\textit{Namen der Teilnehmenden}}
\date{\textit{Datum}}

\begin{document}

\maketitle

\textit{\textbf{Anmerkung:} Alles, was in diesem Dokument kursiv geschrieben ist, muss von der Gruppe entsprechend ihres Themas angepasst (und diese und alle übrigen Anmerkungen entfernt) werden}

\tableofcontents

\newpage

\section{Benutzungsoberfläche}

\textit{\textbf{Anmerkung:} In diesem Abschnitt soll (sofern vorhanden) die Benutzungsoberfläche beschrieben werden; dafür sollen Screenshots verwendet und beschrieben werden, wobei Teile der Oberfläche mit mehreren Komponenten anhand einer Nummerierung des Screenshots (wie im folgenden Beispiel) erklärt werden sollen}

\subsection{Übersicht}

\textit{Kurze Beschreibung der Funktionen der Oberfläche mit (sofern vorhanden) Screenshot der primären Ansicht und Verweis auf die Abbildung \ref{sample-screenshot}}

\begin{figure}[h]
\caption{Beispiel für einen Screenshot mit nummerierten Bereichs-Markierungen}
\label{sample-screenshot}
\includegraphics[width=\textwidth]{sample-screenshot.png}
\end{figure}

\begin{enumerate}
    \item \textit{Beschreibung des mit 1 markierten Teils der Oberfläche}
    \begin{itemize}
        \item \textit{Ggf. stichpunktartige, detailliertere Beschreibung des 1. Teils}
        \item \textit{...}
    \end{itemize}
    \item \textit{...}
\end{enumerate}

\textit{Ggf. weitere Erläuterungen zu der gesamten Oberfläche (falls nötig auch mit weiteren Screenshots)}

\subsection{\textit{Unter-/Detail-Seite 1}}

\textit{Beschreibung der nächsten bzw. einer weiteren Ansicht im gleichen Aufbau wie oben (also mit Screenshot und Erläuterung des Screenshots)}

\subsection{\textit{...}}

\newpage

\section{REST-Schnittstelle}

\textit{\textbf{Anmerkung:} In diesem Abschnitt soll (sofern vorhanden) die REST-Schnittstelle beschrieben werden (als Sicht einer die Software benutzende Person auf das Backend); dafür sollen die einzelnen Endpunkte, die dabei zu verwendende HTTP-Anfragemethode, was sie zurückgeben und (sofern benötigt) die erwarteten Parameter erläutert werden}

\subsection{\textit{Abschnitt 1 der REST-Schnittstelle}}

\textit{Kurze Beschreibung eines ersten Abschnitts der REST-Schnittstelle (z.B. alle Funktion für die An- und Abmeldung an der Oberfläche)\\}

\texttt{\textit{path/to/first/endpoint} $\rightarrow$ \textit{Ergebnis-Format (z.B. \glqq string\grqq)} \dotfill \textit{Methode (z.B. POST)}}
\begin{itemize}[label={}]
    \item \textbf{\textit{nameOfParameter}} (\textit{Parameter-Format (z.B. \glqq string\grqq)}): \textit{Kurze Beschreibung des Parameters}
\end{itemize}

\texttt{\textit{...}}

\subsection{\textit{...}}

\newpage

\section{\textit{Ggf. weitere Schnittstelle(n)}}

\textit{\textbf{Anmerkung:} Sofern die Anwendung über weitere Schnittstellen verfügt (z.B. für die Verwendung über die Kommandozeile) sollen diese hier (und falls nötig in weiteren Abschnitten) beschrieben werden}

\end{document}